\IfFileExists{stacks-project.cls}{%
\documentclass{stacks-project}
}{%
\documentclass{amsart}
}
\usepackage{amssymb}

% For dealing with references we use the comment environment
\usepackage{verbatim}
\usepackage{fancyvrb}
\usepackage{pmboxdraw}
\newenvironment{reference}{\comment}{\endcomment}
%\newenvironment{reference}{}{}
\newenvironment{slogan}{\comment}{\endcomment}
\newenvironment{history}{\comment}{\endcomment}

% For source figures retained as assets
\usepackage{graphicx}

% For commutative diagrams we use Xy-pic
\usepackage[all]{xy}

% We use 2cell for 2-commutative diagrams.
\xyoption{2cell}
\UseAllTwocells

% We use multicol for the list of chapters between chapters
\usepackage{multicol}

% This is generally recommended for better output
\usepackage{lmodern}
\usepackage[T1]{fontenc}
\usepackage{textalpha}

% For cross-file-references
% We prefer the snapshot tmp/xref/#2.aux,
% which latexmkrc replaces atomically, so
% that documents compiled simultaneously
% never read a half-written .aux file.
% The last argument keeps links to #2.pdf.
\usepackage{xr-hyper}
\newcommand{\maybeexternaldocument}[2]{%
  \IfFileExists{#2.aux}{%
    \externaldocument[#1]{#2}%
  }{%
    \PackageWarning{stack}{Missing #2.aux; external references with prefix #1 unavailable}%
  }%
}
\let\externaldocumentorig\externaldocument
\renewcommand{\externaldocument}[2][]{%
  \IfFileExists{tmp/xref/#2.aux}{%
    \externaldocumentorig[#1]{tmp/xref/#2}[#2.pdf]%
  }{%
  \IfFileExists{#2.aux}{%
    \externaldocumentorig[#1]{#2}%
  }{%
    \PackageWarning{stack}{Missing #2.aux; external references with prefix #1 unavailable}%
  }%
  }%
}

% Package for hypertext links:
\usepackage{hyperref}

% For any local file, say "hello.tex" you want to link to please
% use \externaldocument[hello-]{hello}
\externaldocument[introduction-]{introduction}
\externaldocument[conventions-]{conventions}
\externaldocument[sets-]{sets}
\externaldocument[categories-]{categories}
\externaldocument[topology-]{topology}
\externaldocument[sheaves-]{sheaves}
\externaldocument[sites-]{sites}
\externaldocument[stacks-]{stacks}
\externaldocument[fields-]{fields}
\externaldocument[algebra-]{algebra}
\externaldocument[brauer-]{brauer}
\externaldocument[homology-]{homology}
\externaldocument[derived-]{derived}
\externaldocument[simplicial-]{simplicial}
\externaldocument[more-algebra-]{more-algebra}
\externaldocument[smoothing-]{smoothing}
\externaldocument[modules-]{modules}
\externaldocument[sites-modules-]{sites-modules}
\externaldocument[injectives-]{injectives}
\externaldocument[cohomology-]{cohomology}
\externaldocument[sites-cohomology-]{sites-cohomology}
\externaldocument[dga-]{dga}
\externaldocument[dpa-]{dpa}
\externaldocument[sdga-]{sdga}
\externaldocument[hypercovering-]{hypercovering}
\externaldocument[schemes-]{schemes}
\externaldocument[constructions-]{constructions}
\externaldocument[properties-]{properties}
\externaldocument[morphisms-]{morphisms}
\externaldocument[coherent-]{coherent}
\externaldocument[divisors-]{divisors}
\externaldocument[limits-]{limits}
\externaldocument[varieties-]{varieties}
\externaldocument[topologies-]{topologies}
\externaldocument[descent-]{descent}
\externaldocument[perfect-]{perfect}
\externaldocument[more-morphisms-]{more-morphisms}
\externaldocument[flat-]{flat}
\externaldocument[groupoids-]{groupoids}
\externaldocument[more-groupoids-]{more-groupoids}
\externaldocument[etale-]{etale}
\externaldocument[chow-]{chow}
\externaldocument[intersection-]{intersection}
\externaldocument[pic-]{pic}
\externaldocument[weil-]{weil}
\externaldocument[adequate-]{adequate}
\externaldocument[dualizing-]{dualizing}
\externaldocument[duality-]{duality}
\externaldocument[discriminant-]{discriminant}
\externaldocument[derham-]{derham}
\externaldocument[local-cohomology-]{local-cohomology}
\externaldocument[algebraization-]{algebraization}
\externaldocument[curves-]{curves}
\externaldocument[resolve-]{resolve}
\externaldocument[models-]{models}
\externaldocument[functors-]{functors}
\externaldocument[equiv-]{equiv}
\externaldocument[pione-]{pione}
\externaldocument[etale-cohomology-]{etale-cohomology}
\externaldocument[proetale-]{proetale}
\externaldocument[relative-cycles-]{relative-cycles}
\externaldocument[more-etale-]{more-etale}
\externaldocument[trace-]{trace}
\externaldocument[crystalline-]{crystalline}
\externaldocument[spaces-]{spaces}
\externaldocument[spaces-properties-]{spaces-properties}
\externaldocument[spaces-morphisms-]{spaces-morphisms}
\externaldocument[decent-spaces-]{decent-spaces}
\externaldocument[spaces-cohomology-]{spaces-cohomology}
\externaldocument[spaces-limits-]{spaces-limits}
\externaldocument[spaces-divisors-]{spaces-divisors}
\externaldocument[spaces-over-fields-]{spaces-over-fields}
\externaldocument[spaces-topologies-]{spaces-topologies}
\externaldocument[spaces-descent-]{spaces-descent}
\externaldocument[spaces-perfect-]{spaces-perfect}
\externaldocument[spaces-more-morphisms-]{spaces-more-morphisms}
\externaldocument[spaces-flat-]{spaces-flat}
\externaldocument[spaces-groupoids-]{spaces-groupoids}
\externaldocument[spaces-more-groupoids-]{spaces-more-groupoids}
\externaldocument[bootstrap-]{bootstrap}
\externaldocument[spaces-pushouts-]{spaces-pushouts}
\externaldocument[spaces-chow-]{spaces-chow}
\externaldocument[groupoids-quotients-]{groupoids-quotients}
\externaldocument[spaces-more-cohomology-]{spaces-more-cohomology}
\externaldocument[spaces-simplicial-]{spaces-simplicial}
\externaldocument[spaces-duality-]{spaces-duality}
\externaldocument[formal-spaces-]{formal-spaces}
\externaldocument[restricted-]{restricted}
\externaldocument[spaces-resolve-]{spaces-resolve}
\externaldocument[formal-defos-]{formal-defos}
\externaldocument[defos-]{defos}
\externaldocument[cotangent-]{cotangent}
\externaldocument[examples-defos-]{examples-defos}
\externaldocument[algebraic-]{algebraic}
\externaldocument[examples-stacks-]{examples-stacks}
\externaldocument[stacks-sheaves-]{stacks-sheaves}
\externaldocument[criteria-]{criteria}
\externaldocument[artin-]{artin}
\externaldocument[quot-]{quot}
\externaldocument[stacks-properties-]{stacks-properties}
\externaldocument[stacks-morphisms-]{stacks-morphisms}
\externaldocument[stacks-limits-]{stacks-limits}
\externaldocument[stacks-cohomology-]{stacks-cohomology}
\externaldocument[stacks-perfect-]{stacks-perfect}
\externaldocument[stacks-introduction-]{stacks-introduction}
\externaldocument[stacks-more-morphisms-]{stacks-more-morphisms}
\externaldocument[stacks-geometry-]{stacks-geometry}
\externaldocument[moduli-]{moduli}
\externaldocument[moduli-curves-]{moduli-curves}
\externaldocument[examples-]{examples}
\externaldocument[exercises-]{exercises}
\externaldocument[guide-]{guide}
\externaldocument[desirables-]{desirables}
\externaldocument[coding-]{coding}
\externaldocument[obsolete-]{obsolete}
\externaldocument[fdl-]{fdl}
\externaldocument[index-]{index}

% Heap Project documents
\externaldocument[k3-]{k3}
\externaldocument[torelli-]{torelli}
\externaldocument[linear-algebra-]{linear-algebra}
\externaldocument[groups-]{groups}
\externaldocument[complex-analysis-]{complex-analysis}
\externaldocument[functional-analysis-]{functional-analysis}
\externaldocument[categories-]{categories}
\externaldocument[commutative-algebra-]{commutative-algebra}
\externaldocument[ringed-spaces-]{ringed-spaces}
\externaldocument[homological-algebra-]{homological-algebra}
\externaldocument[lie-algebras-]{lie-algebras}
\externaldocument[representation-theory-]{representation-theory}
\externaldocument[smooth-manifolds-]{smooth-manifolds}
\externaldocument[differential-topology-]{differential-topology}
\externaldocument[algebraic-topology-]{algebraic-topology}
\externaldocument[differential-geometry-]{differential-geometry}
\externaldocument[algebraic-geometry-]{algebraic-geometry}
\externaldocument[symplectic-geometry-]{symplectic-geometry}
\externaldocument[complex-geometry-]{complex-geometry}
\externaldocument[hodge-theory-]{hodge-theory}
\externaldocument[%
variations-of-hodge-structures-]{%
variations-of-hodge-structures}
\externaldocument[%
homological-mirror-symmetry-]{%
homological-mirror-symmetry}
\maybeexternaldocument{deformations-}{deformations}
\externaldocument[vertex-algebras-]{vertex-algebras}
\externaldocument[springer-]{springer}
\externaldocument[infty-categories-]{infty-categories}
\externaldocument[seminars-differential-geometry-]{%
seminars-differential-geometry}
\externaldocument[seminars-complex-geometry-]{%
seminars-complex-geometry}
\externaldocument[seminars-algebraic-geometry-]{%
seminars-algebraic-geometry}
\externaldocument[seminars-analysis-pde-]{%
seminars-analysis-pde}
\externaldocument[surfaces-]{surfaces}
\externaldocument[birational-maps-commutative-algebra-]{birational-maps-commutative-algebra}
\externaldocument[cremona-transformations-]{cremona-transformations}
\externaldocument[derived-categories-of-sheaves-]{derived-categories-of-sheaves}
\externaldocument[geometric-invariant-theory-]{geometric-invariant-theory}
\externaldocument[geometric-stability-conditions-and-group-actions-]{geometric-stability-conditions-and-group-actions}
\externaldocument[moduli-spaces-of-sheaves-]{moduli-spaces-of-sheaves}
\externaldocument[bridgeland-stability-general-theory-]{bridgeland-stability-general-theory}
\externaldocument[hilbert-schemes-]{hilbert-schemes}
\externaldocument[noncommutative-abelian-surfaces-and-kummer-type-hyperkahler-varieties-]{noncommutative-abelian-surfaces-and-kummer-type-hyperkahler-varieties}
\externaldocument[mumford-tate-groups-in-hodge-theory-]{mumford-tate-groups-in-hodge-theory}
\externaldocument[hodge-birational-atoms-2026-]{hodge-birational-atoms-2026}
\externaldocument[xxii-egd-2026-]{%
xxii-egd-2026}
\externaldocument[pde-]{pde}
\externaldocument[probability-]{probability}

\externaldocument[linguistics-]{linguistics}

\externaldocument[%
16th-alga-meeting-2026-]{%
16th-alga-meeting-2026}

\externaldocument[%
iv-jornadas-lefschetz-2026-]{%
iv-jornadas-lefschetz-2026}

\externaldocument[reading-the-masters-]{%
reading-the-masters}

\externaldocument[real-analysis-]{%
real-analysis}

\externaldocument[optimal-transport-]{%
optimal-transport}

\externaldocument[history-]{history}

\externaldocument[seiberg-witten-]{seiberg-witten}

\externaldocument[history-philosophy-biology-]{history-philosophy-biology}

\externaldocument[classical-mechanics-]{classical-mechanics}
\externaldocument[electromagnetism-]{electromagnetism}
\externaldocument[waves-and-optics-]{waves-and-optics}
\externaldocument[wif-]{waves-in-fluids-iv}
\externaldocument[thermal-physics-]{thermal-physics}
\externaldocument[quantum-mechanics-]{quantum-mechanics}
\externaldocument[special-relativity-]{special-relativity}
\externaldocument[physics-]{physics}

\externaldocument[chemistry-]{chemistry}

\externaldocument[molecular-cell-biology-]{molecular-cell-biology}

\externaldocument[microbiology-]{microbiology}

\externaldocument[protistology-]{protistology}

\externaldocument[mycology-]{mycology}

\externaldocument[botany-]{botany}

\externaldocument[zoology-]{zoology}

\externaldocument[animal-histology-]{animal-histology}

\externaldocument[animal-physiology-]{animal-physiology}

\externaldocument[systematics-biogeography-]{systematics-biogeography}

\externaldocument[ecology-natural-resources-]{ecology-natural-resources}

\externaldocument[earth-sciences-]{earth-sciences}

\externaldocument[biostatistics-]{biostatistics}

\externaldocument[mathematical-systems-biology-]{mathematical-systems-biology}

\externaldocument[cimpa-el-salvador-2024-]{cimpa-el-salvador-2024}

\externaldocument[anthropology-]{anthropology}

\externaldocument[yang-mills-theory-]{yang-mills-theory}

\externaldocument[seminars-algebra-]{seminars-algebra}

\externaldocument[seminars-physics-]{seminars-physics}

\externaldocument[seminars-computing-]{seminars-computing}

\externaldocument[oktobermat-2026-]{oktobermat-2026}

\externaldocument[saet-2026-]{saet-2026}

\externaldocument[seminars-symplectic-geometry-]{seminars-symplectic-geometry}

\externaldocument[seminars-dynamics-]{seminars-dynamics}

\externaldocument[oktobermat-2024-]{oktobermat-2024}

\externaldocument[quantum-groups-]{quantum-groups}

\externaldocument[legal-2027-]{legal-2027}

\externaldocument[derived-categories-]{derived-categories}

% Theorem environments.
%
\theoremstyle{plain}
\newtheorem{theorem}[subsection]{Theorem}
\newtheorem{proposition}[subsection]{Proposition}
\newtheorem{lemma}[subsection]{Lemma}
\newtheorem{corollary}[subsection]{Corollary}

\theoremstyle{definition}
\newtheorem{definition}[subsection]{Definition}
\newtheorem{example}[subsection]{Example}
\newtheorem{exercise}[subsection]{Exercise}
\newtheorem{situation}[subsection]{Situation}

\theoremstyle{remark}
\newtheorem{remark}[subsection]{Remark}
\newtheorem{remarks}[subsection]{Remarks}

\numberwithin{equation}{subsection}

% Macros
%
\def\lim{\mathop{\mathrm{lim}}\nolimits}
\def\colim{\mathop{\mathrm{colim}}\nolimits}
\def\Spec{\mathop{\mathrm{Spec}}}
\def\Hom{\mathop{\mathrm{Hom}}\nolimits}
\def\Ext{\mathop{\mathrm{Ext}}\nolimits}
\def\SheafHom{\mathop{\mathcal{H}\!\mathit{om}}\nolimits}
\def\SheafExt{\mathop{\mathcal{E}\!\mathit{xt}}\nolimits}
\def\Sch{\mathit{Sch}}
\def\Mor{\mathop{\mathrm{Mor}}\nolimits}
\def\Ob{\mathop{\mathrm{Ob}}\nolimits}
\def\Sh{\mathop{\mathit{Sh}}\nolimits}
\def\NL{\mathop{N\!L}\nolimits}
\def\CH{\mathop{\mathrm{CH}}\nolimits}
\def\proetale{{pro\text{-}\acute{e}tale}}
\def\etale{{\acute{e}tale}}
\def\QCoh{\mathit{QCoh}}
\def\Ker{\mathop{\mathrm{Ker}}}
\def\Im{\mathop{\mathrm{Im}}}
\def\Coker{\mathop{\mathrm{Coker}}}
\def\Coim{\mathop{\mathrm{Coim}}}

% Boxtimes
%
\DeclareMathSymbol{\boxtimes}{\mathbin}{AMSa}{"02}

%
% Macros for moduli stacks/spaces
%
\def\QCohstack{\mathcal{QC}\!\mathit{oh}}
\def\Cohstack{\mathcal{C}\!\mathit{oh}}
\def\Spacesstack{\mathcal{S}\!\mathit{paces}}
\def\Quotfunctor{\mathrm{Quot}}
\def\Hilbfunctor{\mathrm{Hilb}}
\def\Curvesstack{\mathcal{C}\!\mathit{urves}}
\def\Polarizedstack{\mathcal{P}\!\mathit{olarized}}
\def\Complexesstack{\mathcal{C}\!\mathit{omplexes}}
% \Pic is the operator that assigns to X its picard group, usage \Pic(X)
% \Picardstack_{X/B} denotes the Picard stack of X over B
% \Picardfunctor_{X/B} denotes the Picard functor of X over B
\def\Pic{\mathop{\mathrm{Pic}}\nolimits}
\def\Picardstack{\mathcal{P}\!\mathit{ic}}
\def\Picardfunctor{\mathrm{Pic}}
\def\Deformationcategory{\mathcal{D}\!\mathit{ef}}


\begin{document}

\title{Skeletal polytopes}
\maketitle

\phantomsection
\label{section-phantom}

\tableofcontents

\section{Introduction}
\label{section-introduction}

\noindent
In the skeletal approach, a polygon is a
cycle of edges and a polyhedron is a
collection of such cycles. No surface or
solid is enclosed. This chapter records the
basic definitions: skeletal polyhedra and
4-polytopes, abstract polytopes, regularity
and chirality, and realizations.

The text of Sections
\ref{section-historical-background} to
\ref{section-realizations} is taken
verbatim from my master's thesis
\cite{gca-tesina}. Only the typesetting was
changed: cross-references point to the
sections of this chapter, and some conditions
written inside displayed formulas are now
written next to them. The complete thesis,
with its proofs, figures and GAP programs, is
available as the standalone document
\href{https://www.heap-project.org/standalone/chiral-4-polytope.pdf}{A
new chiral 4-polytope}.

\section{Historical background}
\label{section-historical-background}

\noindent
We begin with a brief historical discussion
on the concept of polytope.

We cannot find a moment in history when
triangles and squares came to attract the
attention of humans. Later, when mathematics
became an established discipline, simple
polygons and polyhedra were thoroughly
studied: the first proposition in Euclid's
elements is the construction of a regular
triangle, and the last book is devoted to the
study of Platonic Solids \cite{euclid}.

After the greeks, the definition of polygon
remained essentialy unchanged for many
centuries. The first account of an important
difference dates to the XIV century, when an
Archbishop of Canterbury investigated star
polygons \cite{abstract-polytopes}. These are
essentially different by being non-convex:
their edges intersect in points that are not
vertices of the polygon. They may be,
however, studied by their symmetry properties
just like regular convex polygons.

Regular star polyhedra are the natural
generlization of regular star polygons to
euclidean 3-space. They take us to the XVII
century with Johannes Kepler, who studied the
two whose faces are star polygons
(pentagrams). Their duals, whose faces are
convex but their vertex-figures are not,
where studied in 1809 by Louis Poinsot
\cite{abstract-polytopes}.

Higher-dimensional polytopes were first
studied in the XIX century, when Ludwig
Schläfli found all the regular polytopes
whose symmetry groups are generated by
reflections in hyperplanes in euclidean
spaces \cite{abstract-polytopes}.

Next in history is Coxeter, who, among many
other results, classiffied all discrete
euclidean reflection groups. In the 1930's
him and Petrie redefined the concept of
polytope by letting them have infinite faces,
thus finding three more regular polyhedra in
$\mathbb{E}^3$ \cite{regular-skew}. These
have non-planar vertex-figures.

Fourty years later, in 1977, Grünbaum once
again reintroduced polytopes, this time
permitting the faces to be non-planar
\cite{grunbaum-old-and-new}. The list of
regular polyhedra in $\mathbb{E}^3$ was
extended to 18 in the finite case and 48 in
total \cite{regular-ordinary}. All but one
were classified by Grünbaum, the remaining
one and the completeness of the list are due
to Dress \cite{Dress1981,Dress1985}.

The study of the underlying combinatorial
structure of polytopes led to the concept of
abstract polytope. In 1967 McMullen studied
the lattice of faces of a polytope and
compared its automorphisms with geometric
symmetries \cite{mcmullen-combinatorially}.
In Grümbaum's \cite{grunbaum-regularity}, the
term \textit{polystroma} (stroma=stratum,
layer) is defined to mean an abstract
partially order structure resembling the face
lattice of a polytope.

Chiral polytopes were first studied in this
abstract sense, a general theory first given
in 1991 by Schulte and Weiss
\cite{schulte-chiral}. Back to geometry,
skeletal chiral polyhedra in $\mathbb{E}^3$
were studied in
\cite{chiral-polyhedra-i,chiral-polyhedra-ii},
and the first example of a skeletal chiral
4-polytope in $\mathbb{E}^4$ is Roli's cube
from \cite{rolis-cube}.

\section{Skeletal polyhedra and 4-polytopes}
\label{section-skeletal}

\noindent
Now we give formal definitions and basic
results.

A \textit{skeletal polyhedron} in
$\mathbb{E}^4$ consists of \textit{vertices}
(points in $\mathbb{E}^4$), \textit{edges}
(segments between vertices) and
\textit{faces} (cycles on the graph
determined by the vertices and edges) such
that:
\begin{enumerate}
\item every edge belongs to two faces,
\item the graph determined by the vertices
and edges is connected,
\item every compact subset of $\mathbb{E}^4$
meets finitely many edges, and
\item the \textit{vertex-figure}, defined as
follows, is a connected graph. For any vertex
$v$, the vertices of the vertex-figure are
the neighbours of $v$ and the edges are
segments joining any two neighbours that are
both in some face.
\end{enumerate}
A \textit{skeletal 4-polytope} in
$\mathbb{E}^4$ consits of vertices, edges,
faces and \textit{cells} (skeletal
polyhedra), such that
\begin{enumerate}
\item every face belongs to two cells,
\item the graph determined by the vertices
and the edges is connected,
\item every compact subset of $\mathbb{E}^4$
meets finitely many edges, and
\item the vertex-figure at every vertex is a
skeletal polyhedron.
\end{enumerate}

Hereafter we omit the term "skeletal" when
the context permits it. If $\mathcal{P}$ is a
4-polytope, the set of vertices, edges, faces
and cells is a partial order with respect to
inclusion. We say two such elements are
\textit{incident} if they are comparable.
This poset is called the \textit{abstract
polytope associated to $\mathcal{P}$}. As we
shall see, 4-polytopes defined as above are
realizations of abstract polytopes in the
sense of \cite{abstract-polytopes} and
Section \ref{section-abstract-4-polytopes}.

A \textit{flag} is a 4-tuple of incident
vertex, edge, face and cell. Two flags are
\textit{adjacent} when they differ by only
one element. We call two flags
\textit{0-adjacent} if they differ by a
vertex, \textit{1-adjacent} if they differ by
an edge, \textit{2-adjacent} if they differ
by a face and \textit{3-adjacent} if they
differ by a cell.

A \textit{symmetry} of $\mathcal{P}$ is an
isometry of $\mathbb{E}^4$ that preserves it
set-wise. The group of symmetries of
$\mathcal{P}$ will be denoted by
$G(\mathcal{P})$. We call $\mathcal{P}$
\textit{regular} if $G(\mathcal{P})$ acts
transitively on the set of flags, and
\textit{chiral} if $G(\mathcal{P})$ induces
two orbits on flags so that adjacent flags
are on different orbits. In fact, whether
$\mathcal{P}$ is regular or chiral, its
symmetry group acts transitively on the sets
of vertices, edges, faces and cells.

\section{Abstract 4-polytopes}
\label{section-abstract-4-polytopes}

\noindent
The theory of abstract polytopes is an
essential tool for our work. In this section
we state basic results and constructions upon
which our proofs are based. To avoid using
the word 'abstract' in every definition, we
use the same words as in the previous section
to mean different objects. All polytopes are
abstract up to Section
\ref{section-realizations}, where they are
precisely distinguished.

An \textit{abstract $4$-polytope} is a
partial order $\mathcal{P}$ that satisfies
the following properties (P1) to (P4). The
elements of $\mathcal{P}$ are called
\textit{faces} and the maximally ordered
subsets (chains) are called \textit{flags}.
\begin{enumerate}
\item[(P1)] $\mathcal{P}$ contains a minimal
face and a greatest face, denoted
respectively by $F_{-1}$ and $F_4$.
\item[(P2)] Flags contain $6$ elements,
including $F_{-1}$ and $F_4$.
\end{enumerate}
Since every element in a poset is in some
chain, we may assign to every face a number
from $-1$ to $4$. Such is the \textit{rank
function} of $\mathcal{P}$. We say that two
flags are \textit{$i$-adjacent} if they
differ only by a face of rank $i$.
\begin{enumerate}
\item[(P3)] $\mathcal{P}$ is \textit{strongly
flag connected} in the following sense. For
any two distinct flags $\Phi$ and $\Psi$
there is a sequence
$$
\Phi=\Phi_0,\Phi_1,\ldots,\Phi_k=\Psi
$$
so that $\Phi_{i-1}$ and $\Phi_i$ are
adjacent and $\Phi\cap\Psi\subset\Phi_i$ for
all $i$.
\item[(P4)] If $F$ and $G$ are a $(i-1)$-face
and a $(i+1)$-face with $F<G$ and $0\leq
j\leq2$, then there are exactly two $i$-faces
$H$ such that $F<H<G$.
\end{enumerate}

\section{Regular and chiral abstract
4-polytopes}
\label{section-regular-chiral-abstract}

\noindent
Now we study regular 4-polytopes and their
automorphism groups. An \textit{automorphism}
of $\mathcal{P}$ is a bijection of the
polytope that preserves incidence, and we
denote the automorphism group by
$\Gamma(\mathcal{P})$. $\mathcal{P}$ is
called \textit{regular} if
$\Gamma(\mathcal{P})$ is transitive on flags.

For a regular polytope $\mathcal{P}$ and some
distinguised flag $\Phi$, there exists a
unique involutory automorphism $\rho_j$ of
$\mathcal{P}$ such that $\Phi\rho_i=\Phi^i$
for $i=0,1,2,3$ (Prop. 2B4,
\cite{abstract-polytopes}). These are called
the \textit{distinguished generators} of
$\Gamma(\mathcal{P})$ and satisfy that
\begin{equation}
\label{equation-string-reg}
(\rho_i\rho_j)^{p_{ij}}=\operatorname{Id}
\end{equation}
when $|i-j|\geq2$, for $i,j=0,1,2,3$ (prop.
2B11 \cite{abstract-polytopes}). The numbers
$p_i:=p_{i-1}p_i$ form the \textit{Schläfli
type} $\{p_1,p_2,p_3\}$ of $\mathcal{P}$. The
distinguished generators also satisfy
\begin{equation}
\label{equation-int-reg}
\langle \rho_i|i\in I\rangle\cap
\langle\rho_j|j\in J\rangle=
\langle\rho_i\in I\cap J\rangle
\end{equation}
for $I,J\subseteq\{0,1,2,3\}$, which we call
the \textit{intersection property} (prop.
2B10 \cite{abstract-polytopes}). Any group
$\Gamma$ generated by the involutions
$\rho_0$ to $\rho_3$ satisfying
(\ref{equation-string-reg}) and
(\ref{equation-int-reg}) is called a
\textit{C-string group}. We have established
that $\Gamma(\mathcal{P})$ is one such group.

Conversely, we may construct a regular
abstract $4$-polytope from a C-string group
$\Gamma$. Define
$\Gamma_j:=\langle\rho_j|i\neq j\rangle$ and
$\Gamma_{-1}=\Gamma_4=\Gamma$, and take the
set of $i$-faces to be the set of right
cosets $\Gamma_i\varphi$ for
$\varphi\in\Gamma$. Then there is an order
relation with respect to which the set of all
faces is a regular $4$-polytope whose
automorphism group is $\Gamma$ (thm. 2E11,
\cite{abstract-polytopes}).

\medskip\noindent
Now we revise the analogue construction for
chiral polytopes. An abstract 4-polytope
$\mathcal{P}$ is called \textit{chiral} if
$\Gamma(\mathcal{P})$ induces two orbits on
flags and two adjacent flags are in different
orbits.

First observe that by defining
$\sigma_i:=\rho_{i-1}\rho_i$ for the
distinguished generators of a regular
polytope with bipartite flags, it holds that
\begin{equation}
\label{equation-chiral-rels}
\begin{aligned}
\sigma_1^{p_1}=\sigma_2^{p_2}
&=\sigma_3^{p_3}=\operatorname{Id}\\
(\sigma_1\sigma_2)^2
=(\sigma_1\sigma_2&\sigma_3)^2
=(\sigma_2\sigma_3)^2=\operatorname{Id}
\end{aligned}
\end{equation}
for some positive integers $p_1$, $p_2$ and
$p_3$ (see \cite{schulte-chiral}).

The group generated by these elements is
called the \textit{rotation subgroup} of
$\mathcal{P}$, denoted by
$\Gamma^+(\mathcal{P})$. For chiral polytopes
this will be the whole group: if
$\mathcal{P}$ is chiral,
$\Gamma(\mathcal{P})$ is generated by three
automorphisms $\sigma_1,\sigma_2$ and
$\sigma_3$ that satisfy
(\ref{equation-chiral-rels}) (prop. 3,
\cite{schulte-chiral}).

To construct a chiral or regular abstract
polytope from a group we must also require
some sort of intersection property. It will
suffice that
\begin{equation}
\label{equation-chiral-int}
\begin{aligned}
\langle\sigma_1\rangle\cap
\langle\sigma_2\rangle=\{\operatorname{Id}\}
&=\langle\sigma_2\rangle\cap
\langle\sigma_3\rangle\\
\langle\sigma_1,\sigma_2\rangle\cap
\langle\sigma_2,&\sigma_3\rangle
=\langle\sigma_2\rangle.
\end{aligned}
\end{equation}
Let $\Gamma=\langle
\sigma_1,\sigma_2,\sigma_3\rangle$ be a group
satisfying both (\ref{equation-chiral-rels})
and (\ref{equation-chiral-int}). Define
$$
\begin{aligned}
\Gamma^0&=\langle\sigma_2,\sigma_3\rangle,\\
\Gamma^1&=\langle \sigma_3,\sigma_1\sigma_2\rangle,\\
\Gamma^2&=\langle \sigma_1,\sigma_2\sigma_3\rangle,\\
\Gamma^3&=\langle \sigma_1,\sigma_2\rangle,
\end{aligned}
$$
and $\Gamma_{-1}=\Gamma^4=\Gamma$, and take
the $i$-faces to be the right cosets
$\Gamma_i\varphi$ for $\varphi\in\Gamma$.
Again, the set of all faces admits a partial
order with respect to which it is a chiral or
a regular abstract 4-polytope $\mathcal{P}$
such that $\Gamma^+(\mathcal{P})=\Gamma$
(thm. 1 and lem. 11, \cite{schulte-chiral}).

In this construction we may distinguish
chiral from regular abstract 4-polytopes as
follows. $\Gamma^+(\mathcal{P})$ is of index
2 in $\Gamma(\mathcal{P})$ if and only if
there exists an involutory automorphism
$\rho:\Gamma\to \Gamma$ such that
\begin{equation}
\label{equation-combinatorially-chiral}
\rho(\sigma_1)=\sigma_1^{-1},\quad
\rho(\sigma_2)=\sigma_1^2\sigma_2,\quad
\rho(\sigma_3)=\sigma_3,
\end{equation}
in which case $\mathcal{P}$ cannot be chiral
(thm. 1, \cite{schulte-chiral}).

\section{Realizations}
\label{section-realizations}

\noindent
Now we review the relationship between
abstract and skeletal polytopes as defined in
Section \ref{section-skeletal}.

A \textit{realization} of an abstract
polytope $\mathcal{P}$ is a map $\beta$ from
the set of abstract 0-faces $\mathcal{P}_0$
into $\mathbb{E}^4$, so that the set
$V_0:=\mathcal{P}_0\beta$ is the set of
geometric vertices. All other geometric faces
are defined by functions from the set of
abstract $i$-faces $\mathcal{P}_i$ to some
nested power set of the geometric vertex set:
edges are sets of vertices, faces are sets of
edges and so on.

Formally, for $i=1,2,3$, $\beta$ induces a
surjection $\beta_i:\mathcal{P}_i\to V_i$,
where $V_i$ is thought as the set of
geometric $i$-faces, consisting of the
elements
$$
F\beta_i:=\{G\beta_{i-1}|
G\in\mathcal{P}_{i-1},\ G\leq F\}
$$
for $F\in\mathcal{P}_i$ (thm. 5A1,
\cite{abstract-polytopes}). For example, an
abstract edge $E\in\mathcal{P}_1$ is mapped
to a set consisting of the two points
(there's only two by (P4)) in $\mathbb{E}^4$
whose preimages are abstract vertices smaller
than $E$ in $\mathcal{P}$.

Of course, we expect the number of $i$-faces
of the abstract and geometric structures to
be the same. A realization is
\textit{faithful} if every $\beta_i$ is a
bijection.

We also expect automorphisms to correspond
with isometries of $\mathbb{E}^4$. A
realization is \textit{symmetric} when every
automorphism of $\mathcal{P}$ induces a
permutation of $V_0$, which in turn
determines a unique isometry of
$\mathbb{E}^4$ (if the vertex set affinely
spans $\mathbb{E}^4$). Then these isometries
are an euclidean representation of
$\Gamma(\mathcal{P})$.

Conversely, given an euclidean representation
of the automorphism group
$\Gamma(\mathcal{P})$ of an abstract
4-polytope, we obtain a realization by
\textit{Wythoff's construction} as follows.

For the regular case, let $\langle R_0,
R_1,R_2,R_3\rangle$ be an euclidean
representation of the automorphism group of a
regular abstract 4-polytope $\mathcal{P}$.
Should there be any, define a point
$v\in\mathbb{E}^4$ that is fixed by $R_1$,
$R_2$ and $R_3$ but not by $R_0$ as the base
vertex and take the orbit of $v$ for the
geometric vertex-set of a realization.
Further, define the base edge $e=v\langle
R_0\rangle$, the base face $f=e\langle
R_0,R_1\rangle$ and the base cell $c=f\langle
R_0,R_1,R_2\rangle$ for a geometric base
flag.

Now let $\langle S_1,S_2,S_3\rangle$ be a
representation of the automorphism group, or
rotation subgroup, of a chiral or regular
abstract 4-polytope, respectively. For a
geometric base vertex choose any point that
is fixed by $S_2$ and $S_3$ but not by $S_1$.
The orbit of this point is the geometric
vertex-set of a realization. For a geometric
flag we define the base edge as $e=v\langle
S_1S_2\rangle$, or equivalently, as
$e=\{v,vS_1^{-1}\}$ since $S_1^{-1}=R_1R_0$
when $\mathcal{P}$ is regular. Define the
base face as $f=e\langle S_1\rangle$ and the
base cell as $c=f\langle S_1,S_2\rangle$.

While it is true that realizations of
4-polytopes are mere clouds of points, they
are no generalization of skeletal
polytopes---the segment between vertices is
still taken to be flat. Save this technical
difference, it is straightforward to check
that a regular or chiral skeletal 4-polytope
as defined in Section \ref{section-skeletal}
is a faithful and symmetric realization of
the abstract 4-polytope related to it.

\medskip\noindent
It remains only to formally define star
polytopes. Following Section 7D from
\cite{abstract-polytopes}, we say a
faithfully realized abstract 4-polytope is
\textit{classical} if every geometric
$i$-face has dimension $i$ (its affine hull
is $i$-dimensional). For any such geometric
regular polytope, we may write its generating
hyperplane reflections as
$$
R_i=\{x\in\mathbb{E}^4|
\langle x,u_i\rangle=0\}
$$
for some unit vectors $u_i$. Then $\langle
u_i,u_j\rangle=0$ for $|i-j|\geq2$, and,
since $\langle R_i,R_j\rangle$ is a finite
group, there is a rational number $p_i$ such
that $\langle
u_{i-1},u_i\rangle=\cos(\pi/p_i)$. We define
the \textit{Shläfli type} of this polytope as
$\{p_1,p_2,p_3\}$ and call it a \textit{star
4-polytope} if any of the $p_i$ is not an
integer. Of course, the definition for
polyhedra is analogous.

\bibliography{my}
\bibliographystyle{amsalpha}

\end{document}
